\chapter{Первые шаги}
\label{ch:getting-started}

\section{Регистрация}
\label{sec:getting-started:register}

Откройте главную страницу сайта и нажмите ``Регистрация'' в правом верхнем углу --- это страница \texttt{/login?register=1}. Понадобятся три поля: адрес почты, отображаемое имя и пароль. Почта служит логином, имя видят остальные участники игры.

\screenshotOptional{user/login.png}{страница входа и регистрации}

После регистрации вы сразу вошли в систему. Аккаунт создаётся как обычный игрок: вы можете участвовать в играх, но не можете создавать сценарии. Как получить права мастера --- в разделе \ref{sec:getting-started:master-rights}.

\section{Вход}
\label{sec:getting-started:login}

Войти можно тремя способами.

\paragraph{Почта и пароль}
Обычный вход на \texttt{/login}.

\paragraph{Ссылка для входа (magic link)}
Одноразовая ссылка вида \texttt{/auth/link/<токен>}. Её выдаёт Telegram-бот (см.\ главу \ref{ch:telegram}). Переход по ссылке авторизует вас без ввода пароля; повторно та же ссылка не сработает.

\paragraph{Telegram}
Если вы открыли сайт кнопкой внутри Telegram, вход происходит автоматически --- ничего вводить не нужно.

\begin{quote}
\small
Обратите внимание: кнопки ``Войти через Telegram'' на странице входа нет. В обычном браузере этот способ недоступен --- нужна либо ссылка от бота, либо открытие сайта прямо из Telegram.
\end{quote}

\paragraph{Сколько держится вход}
Сессия живёт несколько дней и продлевается сама, пока вы пользуетесь сайтом. Если вход всё же слетел, вас перебросит на \texttt{/login}, а после входа вернёт на ту страницу, с которой выкинуло.

\section{Профиль}
\label{sec:getting-started:profile}

Страница \texttt{/me} --- личный кабинет.

\screenshotOptional{user/profile.png}{личный кабинет}

Здесь:

\begin{itemize}
\item имя, почта и аватар, кнопка ``Редактировать'';
\item блок ``Мастер'' с заявками --- он появляется только у мастеров, и на кнопке видно число заявок, ждущих ответа;
\item ``Запущенные сценарии'' --- миры, которые вы ведёте, со ссылками в мир и в исходный сценарий;
\item ``Броски и статистика'' --- сводка и переход ко всем броскам;
\item ``История сессий'' --- игры, в которых вы участвовали, с пометкой ``активна'' у идущих и ссылкой ``Перейти в сессию''.
\end{itemize}

Отдельная страница \texttt{/me/rolls} хранит историю всех ваших бросков --- по ней удобно восстановить, что и когда выпало, если возник спор за столом.

\section{Права мастера}
\label{sec:getting-started:master-rights}

Создавать сценарии и вести игры может не каждый: нужен признак ``мастер''. Сам себе его выдать нельзя --- это делает администратор платформы. Пока признака нет, разделы ``Сценарии'', ``Кампании'' и ``Запущенные'' будут открываться, но окажутся пустыми: создавать в них нечего.

Если вы разворачиваете платформу для себя, первому пользователю права выдаются напрямую в базе --- см.\ раздел \ref{sec:admin:grant-master}.

\section{Что доступно без регистрации}
\label{sec:getting-started:anonymous}

Не заводя аккаунт, можно смотреть чужую игру в режиме наблюдателя: мастер присылает ссылку на комнату, вы её открываете и видите происходящее. Подробности --- в главе \ref{ch:observer}.
